% typo_francaise.tex
% Règles typographiques françaises (guide du dactylographe).
% Inclusion : % typo_francaise.tex
% Règles typographiques françaises (guide du dactylographe).
% Inclusion : % typo_francaise.tex
% Règles typographiques françaises (guide du dactylographe).
% Inclusion : % typo_francaise.tex
% Règles typographiques françaises (guide du dactylographe).
% Inclusion : \input{typo_francaise} dans le préambule.
% Pré-requis : babel chargé avec l'option [french] (gère déjà les espaces
% fines insécables avant : ; ! ? et autour de « »).

% --- Espace fine insécable automatique avant : ; ! ? -------------------
% Force babel-french à insérer une fine insécable avant la ponctuation double
% quelle que soit la saisie (mot:, mot :, mot   : donnent le même rendu).
\frenchsetup{AutoSpacePunctuation=true}

% --- Guillemets français : taper "texte" produit « texte » ---------------
\usepackage{csquotes}
\MakeOuterQuote{"}

% --- Nombres : 1234567 -> 1 234 567 ; décimales avec virgule -------------
% Usage : \num{1234567}, \num{3,14}, \SI{12}{\kilo\meter}, \SI{50}{\percent}
\usepackage{siunitx}
\sisetup{
  locale                  = FR,
  group-separator         = {\,},   % espace fine entre milliers
  group-minimum-digits    = 4,      % regrouper à partir de 4 chiffres
  output-decimal-marker   = {,},
  per-mode                = symbol,
  detect-all                        % suit la fonte du texte courant
}

% --- Micro-typographie : kerning, justification, débord en marge ---------
\usepackage{microtype}

% --- Tirets ---------------------------------------------------------------
% Sous LaTeX : --- = tiret cadratin (—), -- = demi-cadratin (–), - = trait d'union.
% Aucune configuration nécessaire, juste un rappel d'usage.

% --- Note d'usage ---------------------------------------------------------
% Espaces insécables manuelles (toujours nécessaires hors ponctuation forte) :
%   M.~Dupont, 12~km, n\textsuperscript{o}~3, p.~42, etc.
% Les espaces fines insécables avant : ; ! ? sont automatiques (babel french).
 dans le préambule.
% Pré-requis : babel chargé avec l'option [french] (gère déjà les espaces
% fines insécables avant : ; ! ? et autour de « »).

% --- Espace fine insécable automatique avant : ; ! ? -------------------
% Force babel-french à insérer une fine insécable avant la ponctuation double
% quelle que soit la saisie (mot:, mot :, mot   : donnent le même rendu).
\frenchsetup{AutoSpacePunctuation=true}

% --- Guillemets français : taper "texte" produit « texte » ---------------
\usepackage{csquotes}
\MakeOuterQuote{"}

% --- Nombres : 1234567 -> 1 234 567 ; décimales avec virgule -------------
% Usage : \num{1234567}, \num{3,14}, \SI{12}{\kilo\meter}, \SI{50}{\percent}
\usepackage{siunitx}
\sisetup{
  locale                  = FR,
  group-separator         = {\,},   % espace fine entre milliers
  group-minimum-digits    = 4,      % regrouper à partir de 4 chiffres
  output-decimal-marker   = {,},
  per-mode                = symbol,
  detect-all                        % suit la fonte du texte courant
}

% --- Micro-typographie : kerning, justification, débord en marge ---------
\usepackage{microtype}

% --- Tirets ---------------------------------------------------------------
% Sous LaTeX : --- = tiret cadratin (—), -- = demi-cadratin (–), - = trait d'union.
% Aucune configuration nécessaire, juste un rappel d'usage.

% --- Note d'usage ---------------------------------------------------------
% Espaces insécables manuelles (toujours nécessaires hors ponctuation forte) :
%   M.~Dupont, 12~km, n\textsuperscript{o}~3, p.~42, etc.
% Les espaces fines insécables avant : ; ! ? sont automatiques (babel french).
 dans le préambule.
% Pré-requis : babel chargé avec l'option [french] (gère déjà les espaces
% fines insécables avant : ; ! ? et autour de « »).

% --- Espace fine insécable automatique avant : ; ! ? -------------------
% Force babel-french à insérer une fine insécable avant la ponctuation double
% quelle que soit la saisie (mot:, mot :, mot   : donnent le même rendu).
\frenchsetup{AutoSpacePunctuation=true}

% --- Guillemets français : taper "texte" produit « texte » ---------------
\usepackage{csquotes}
\MakeOuterQuote{"}

% --- Nombres : 1234567 -> 1 234 567 ; décimales avec virgule -------------
% Usage : \num{1234567}, \num{3,14}, \SI{12}{\kilo\meter}, \SI{50}{\percent}
\usepackage{siunitx}
\sisetup{
  locale                  = FR,
  group-separator         = {\,},   % espace fine entre milliers
  group-minimum-digits    = 4,      % regrouper à partir de 4 chiffres
  output-decimal-marker   = {,},
  per-mode                = symbol,
  detect-all                        % suit la fonte du texte courant
}

% --- Micro-typographie : kerning, justification, débord en marge ---------
\usepackage{microtype}

% --- Tirets ---------------------------------------------------------------
% Sous LaTeX : --- = tiret cadratin (—), -- = demi-cadratin (–), - = trait d'union.
% Aucune configuration nécessaire, juste un rappel d'usage.

% --- Note d'usage ---------------------------------------------------------
% Espaces insécables manuelles (toujours nécessaires hors ponctuation forte) :
%   M.~Dupont, 12~km, n\textsuperscript{o}~3, p.~42, etc.
% Les espaces fines insécables avant : ; ! ? sont automatiques (babel french).
 dans le préambule.
% Pré-requis : babel chargé avec l'option [french] (gère déjà les espaces
% fines insécables avant : ; ! ? et autour de « »).

% --- Espace fine insécable automatique avant : ; ! ? -------------------
% Force babel-french à insérer une fine insécable avant la ponctuation double
% quelle que soit la saisie (mot:, mot :, mot   : donnent le même rendu).
\frenchsetup{AutoSpacePunctuation=true}

% --- Guillemets français : taper "texte" produit « texte » ---------------
\usepackage{csquotes}
\MakeOuterQuote{"}

% --- Nombres : 1234567 -> 1 234 567 ; décimales avec virgule -------------
% Usage : \num{1234567}, \num{3,14}, \SI{12}{\kilo\meter}, \SI{50}{\percent}
\usepackage{siunitx}
\sisetup{
  locale                  = FR,
  group-separator         = {\,},   % espace fine entre milliers
  group-minimum-digits    = 4,      % regrouper à partir de 4 chiffres
  output-decimal-marker   = {,},
  per-mode                = symbol,
  detect-all                        % suit la fonte du texte courant
}

% --- Micro-typographie : kerning, justification, débord en marge ---------
\usepackage{microtype}

% --- Tirets ---------------------------------------------------------------
% Sous LaTeX : --- = tiret cadratin (—), -- = demi-cadratin (–), - = trait d'union.
% Aucune configuration nécessaire, juste un rappel d'usage.

% --- Note d'usage ---------------------------------------------------------
% Espaces insécables manuelles (toujours nécessaires hors ponctuation forte) :
%   M.~Dupont, 12~km, n\textsuperscript{o}~3, p.~42, etc.
% Les espaces fines insécables avant : ; ! ? sont automatiques (babel french).
